% 原文 p.85（後半）


\subsectionlink{環境への情報の冗長な符号化と「量子Darwin主義」}
これまでは、系の水準で何を観測できるかに注目してきた。環境は、系についての経路情報、より一般的には、どの状態かという情報を運び去る「排出先」としてだけ働いていた。環境をさらに観測したり、相互作用したりしないと仮定した。同時に、系と環境の相互作用が、系についての情報を絶えず環境に符号化するので、環境は、系の情報を間接的に獲得するための、巨大な資源を構成する。実際、すべてでないにせよ大半の場合、観測者が系の状態について情報を集めるのは、間接的な観測によることを認識するのが重要である。すなわち、過去に系と相互作用したため、その状態についての情報を含む環境自由度の断片を、捕らえるのである。そのような間接的情報獲得の、おそらく最も普通の例は、関心のある物体に散乱された光子を、視覚的に記録することである。私たちは、物と直接相互作用して見るのでなく、その空間構造を「測定」した光を通じて見る。

なぜ、環境自由度を捕らえて系を間接観測することの重要性という、この認識が、それほど大切なのだろうか。これに答えるため、2.1.1節で、古典物理学の特徴として、初めは状態をまったく知らなくてもよい多数の独立した観測者が、系の状態を乱さずに、それを知り、互いに一致できることを指摘したのを思い出そう。この意味で、古典状態は客観的にあらかじめ存在し、「古典的実在」という概念をもたらす。これに対し、閉じた量子系への測定は、一般に系の状態を変える。

% 原文 p.86
\hypertarget{srcpage:86}{}\label{srcpage:86}
したがって、閉じた系の量子状態を、古典状態と同じ仕方で存在するものと見なすことは不可能である。そこで、量子的基盤の内部から、古典的実在がどう現れるのか、すなわち、上の意味で観測可能量がどう「客観化」されるかという問いが生じる。前の2.8節で論じた優先状態の環境誘起超選択は、系のHilbert空間の可能な状態の、ある部分集合だけが実際に観測される理由を説明し、この問いへの答えに、確かに大きく寄与した。それでも、系への直接測定はどれも、一般になお系の状態を変えるので、前段落で述べた問題は残る。

Los AlamosのZurek、Ollivier、Poulin、Blume-Kohoutらのグループ\cite{ref016,ref066,ref076,ref084,ref104,ref110,ref111,ref112,ref113}が、初めて明瞭に認識して系統的に調べたように、どの状態かを監視する環境の働きが、ここでもう一つの重要な役割を果たす。系の情報は環境に符号化されるので、系自体と直接相互作用して乱すことなく、それを獲得できる。こうして環境の役割は、関心のある系の優先状態の選択と、それらの状態間の局所的位相コヒーレンスの非局在化から、系の状態に関する情報の伝達へと広がる。そこで、Los Alamosのグループが最初に明示した中心的な問いは、次となる。どのような種類の情報が、どのようにして、環境の多数の異なる断片に冗長かつ頑健に記憶され、多数の観測者が系の状態を乱さずに取り出せ、その状態の実効的な古典性を達成するのか\footnote{量子測定における冗長性の重要性は、すでに1980年代初頭にZurek\cite{ref076}が初めて強調した（文献\cite{ref009}も参照）。}。

この問いに答えるLos Alamosのグループの研究は、「証人としての環境」という研究計画、すなわち通信路としての環境の役割の認識と、量子Darwin主義、すなわち系についてのどの情報を、環境が安定に記憶し増殖させられるかという研究の名のもとで行われている\cite{ref066,ref110,ref111,ref112,ref113}。環境に刻まれた情報の完全性と冗長性の程度を明示的に定量化するためには、通常、古典的\cite{ref110,ref111}または量子的\cite{ref016,ref112,ref113}な相互情報量という尺度を使う。大まかに言えば、これは、環境またはその断片を測定して獲得できる、系についての情報量を表し、Shannonエントロピー\cite{ref114,ref115,ref116}またはvon Neumannエントロピーで表現される\footnote{各断片に含まれる情報量は、系自体が提供する最大情報量、すなわち系のvon Neumannエントロピーより、常にいくらか小さいことに注意しよう\cite{ref113}。}。

% 原文 p.87
\hypertarget{srcpage:87}{}\label{srcpage:87}
古典的相互情報量という尺度は、系$\mathcal S$と環境$\mathcal E$の特定の観測可能量の選択に基づき、$\mathcal E$の一部の何らかの観測可能量を測定することによって、与えられた$\mathcal S$の観測可能量の測定結果を、どれほどよく予言できるかを定量化する\cite{ref110,ref111}。Blume-KohoutとZurekの、より最近の研究\cite{ref112,ref113}で使われる量子相互情報量$I_{\mathcal S:\mathcal E}$は、古典的相互情報量の一般化と見なせ、$I_{\mathcal S:\mathcal E}\equiv S(\hat\rho_{\mathcal S})+S(\hat\rho_{\mathcal E})-S(\hat\rho)$と定義される。ここで$\hat\rho_{\mathcal S}$、$\hat\rho_{\mathcal E}$、$\hat\rho$は、それぞれ$\mathcal S$、$\mathcal E$、複合系$\mathcal{SE}$の密度行列であり、$S(\hat\rho)$はvon Neumannエントロピー(2.32)である。したがって$I_{\mathcal S:\mathcal E}$は、$\mathcal S$、$\mathcal E$間のすべての量子的相関を破壊して生じるエントロピー量を測り、すなわち両者の相関の程度を定量化する。部分系のShannonエントロピーに基づく古典的相互情報量と、von Neumannエントロピーで定義する量子相互情報量という、二つの尺度から導かれる結果は、ほぼ一致することが見いだされている\cite{ref016,ref110,ref111,ref112,ref113}。具体的な尺度の選択に対するこの頑健性は、OllivierとZurek\cite{ref117}が示したように、二尺度の差、いわゆる量子不協和\cite{ref117}が、デコヒーレンスが十分有効で、明確な指針基底が選ばれると消失するためである。したがって、捕らえる断片以外の環境全体によってデコヒーレンスした系についての情報の、量子不協和は小さくなければならない。

Ollivier、Poulin、Zurek\cite{ref110,ref111}、およびBlume-KohoutとZurek\cite{ref112,ref113}は、環境の多数の異なる部分集合に、最も完全かつ冗長に刻める系の観測可能量は、2.8節で論じた系と環境の相互作用が選ぶ指針観測可能量に一致することを見いだした。逆に、他の大半の状態は、冗長には記憶できないように見える。したがって、系の指針状態の最大情報のうち、かなりの量を推論するには、環境のごく小さな割合を調べれば足りる。他方、同じ断片で他の観測可能量を測ろうとすれば、対応する系の観測可能量の情報は冗長に記憶されないため、ほとんど何も分からない。したがって、環境が超選択する系の状態は、二重の役割を持つ。環境との相互作用に最も乱されない状態であり、環境自由度を調べることにより、系を乱さずに、最も容易に見いだせる状態でもある。指針観測可能量についての同じ情報が、環境の多数の断片に独立に記憶されるので、多数の観測者が異なる断片でこの観測可能量を測定しても、知見は自動的に一致する。この意味で、環境が超選択する状態に、実効的な客観的存在を帰属させられる（8.2.3節も参照）。

% 原文 p.88
\hypertarget{srcpage:88}{}\label{srcpage:88}
本節の方針に沿って観測可能量の客観化を調べるLos Alamosグループの研究は、まだ始まったばかりである。客観化過程の具体的な力学的発展\cite{ref111}や、「冗長性による客観化」という考えの現在の扱いにおける、仮定と定義の役割などの重要な側面は、より詳細で現実的な系と環境の模型を含む研究とともに、なお調査中である。しかし、閉じた系という見方を離れ、環境に冗長かつ頑健に記憶された情報を捕らえることとして観測を記述する方法が、量子的領域からの古典性の出現を、純粋に量子力学的に説明する、きわめて有望な候補であることは、すでに明らかになったはずである。
